\documentclass[11pt]{article}
\usepackage[margin=.85in]{geometry}
\usepackage{amsmath,amssymb}
\usepackage[colorlinks=true,urlcolor=blue]{hyperref}
\setlength{\parindent}{0pt}
\setlength{\parskip}{.42em}
\setlength{\emergencystretch}{2em}
\newcommand{\ch}{\chi_p}
\newcommand{\Q}{\mathbb Q}
\newcommand{\Z}{\mathbb Z}
\newcommand{\src}{\href{https://www.kurims.kyoto-u.ac.jp/EMIS/journals/AMUC/_vol-71/_no_2/_mollin/mollin.pdf}{Mollin--Cheng--Goddard, Theorems 2.1--2.2 and Remark 2.1}}
\newcommand{\cmqr}{\href{https://arxiv.org/pdf/2404.08380}{Carneiro--Milinovich--Quesada-Herrera--Ramos, proof of Theorem 6}}
\begin{document}
{\Large\bfseries Which nonresidue primes occur in prescribed ideal cycles?}\\[2pt]
{\large A fixed-prime estimate for \(p\equiv23\pmod{24}\)}

Put \(L=\log p\), \(B=\eta L^2\) with fixed \(\eta>0\), and \(\ch(n)=(n/p)\).
We take \(p\) large enough that \(5\le B<\sqrt p\) and \(2\sqrt p<p\).
We seek a nonresidue at most \(B\) for every sufficiently large \(p\).
First test all primes \(s\le B\) with \(s\equiv1\bmod4\).
An inert \(s\) is itself a witness by reciprocity. In the remaining case
all these primes split in \(K=\Q(\sqrt p)\). For example, the first test
already detects 41 at \(p=2{,}080{,}679\), since \(p\equiv11\bmod41\)
and \(11^{20}\equiv-1\bmod41\).

\section*{1. The exact average over the prescribed starts}
The full ring of integers is \(O=\Z[\sqrt p]\). Its discriminant is \(4p\).
Each reduced primitive ideal has a unique description
\(I=(Q,r+\sqrt p)\) with \(Q\mid p-r^2\), \(\sqrt p-Q<r<\sqrt p\),
and \(Q<r+\sqrt p\). Its norm is \(Q<2\sqrt p\). Every primitive ideal of
norm \(<\sqrt p\) is reduced. A complete period consists of all the reduced
ideals of its \emph{ordinary} ideal class. In particular, each of the two
roots mod \(s\) gives the prescribed reduced starting ideal. \src.

From a state \((r,Q)\) the next state \((a,Q')\) is
\[
 a=\left\lfloor(r+\sqrt p)/Q\right\rfloor Q-r,\qquad Q'=(p-a^2)/Q.
\]
Thus \(\ch(Q')=-\ch(Q)\), since \(a^2\equiv-QQ'\bmod p\).
Every odd prime \(q\mid Q'\) splits in \(K\), and
\(\ch(q)=(-1/q)\). Hence a first neighbor of any start contains a prime
\(q\equiv3\bmod4\) to odd exponent. Two is a residue.

Let \(G=\operatorname{Cl}(O)\), \(h=|G|\), and \(w_C\) count the prescribed
starts in \(C\). Let \(S=\{C:w_C>0\}\), \(T=\{q\le B:q\text{ prime},
\ q\equiv3\bmod4,\ (p/q)=1\}\), and define
\[
 \nu_C=\sum_{\substack{I\ {\rm reduced}\\ [I]=C}}\sum_{q\in T}
             {\bf1}_{q\mid NI}.
\]
For \(w=\sum_Cw_C\), the requested mean of the numbers of incidences is exactly
\begin{equation}
 J_p(B)=\frac1w\sum_{C\in G}w_C\nu_C .
\end{equation}
An incidence counts a pair consisting of an entry and a prime divisor, once
regardless of the exponent. For \(q\in T\) both prime ideals of norm \(q\)
occur in the cycles of their own classes, since \(q<\sqrt p\). Composite norms in other classes
can also be divisible by \(q\). Therefore full class coverage would detect
all of \(T\), but even full coverage needs \(T\ne\varnothing\). In the
remaining case this local existence statement is precisely \(n_p\le B\).

\newpage
\section*{2. A first-step count without an assumed small nonresidue}
Let \(W_1(x)\) be the set of distinct primes \(q\equiv3\bmod4\), \(q\le x\),
dividing a first neighbor of some prescribed start. The following bound uses
only elementary arithmetic at the given \(p\):
\begin{equation}
 \sum_{q\in W_1(2\sqrt p)}\left\lceil B/q\right\rceil
       \log(2B\sqrt p/q)\ \ge\
 \vartheta_1(B),\qquad
 \vartheta_1(B)=\sum_{\substack{s\le B\ {\rm prime}\\ s\equiv1\bmod4}}\log s .
\end{equation}
In particular one has the disjunction
\begin{equation}
 W_1(B)\ne\varnothing\qquad\text{or}\qquad
 \#\{q\in W_1(2\sqrt p):q>B\}
 \ \ge\ \frac{\vartheta_1(B)}{\log(2\sqrt p)}
 = (1+o_\eta(1))\frac{B}{\log p}.
\end{equation}
The second alternative supplies order \(\eta\log p\) \emph{distinct}
nonresidue primes. The upper endpoint for their sizes is \(2\sqrt p\).

To prove this, group the starts by the integer \(a\) in the next state.
For fixed \(s\), this \(a\) is the representative of the opposite root
in \((\sqrt p-s,\sqrt p)\), so each of the two possible \(a\)'s occurs once.
For an occupied \(a\), let \(D_a\) be the product of the distinct starting
primes giving that \(a\). Then \(D_a\mid p-a^2\),
\(\ch(D_a)=1\), and
\[
 \sum_a\log D_a=2\vartheta_1(B),\qquad
 \ch\big((p-a^2)/D_a\big)=-1 .
\]
Choose a nonresidue prime divisor \(q_a\) of that quotient. It divides
the first neighbor for every start with this \(a\), hence lies in
\(W_1(2\sqrt p)\). Moreover
\(q_aD_a\le p-a^2<2B\sqrt p\). For each fixed \(q\), all \(a\)'s so
assigned are in two residue classes modulo \(q\) in an interval of length
\(B\). There are at most \(2\lceil B/q\rceil\) of them. Summing the
resulting bounds for \(\log D_a\) proves (2). If all assigned primes exceed
\(B\), each occurs at most twice and costs at most \(\log(2\sqrt p)\) per
occurrence. This gives (3), using the PNT in the fixed progression \(1\bmod4\).

The same argument gives the computable size bound
\[
 \min W_1(2\sqrt p)\le \min_{a:D_a>1}(p-a^2)/D_a.
\]
A sufficiently large product \(D_a\) at one \(a\) would bring the right side
below \(B\). The sum of its logarithms fixes only the total mass.
The supply of large primes allows the second alternative in (3).
Dirichlet GRH gives \(\asymp\sqrt p/\log p\) split primes \(3\bmod4\)
in \((\sqrt p,2\sqrt p)\), versus \(\asymp_\eta\log p\) required by (3).
Indeed the progression and splitting indicators use only the four
characters \(1,\chi_{-4},\ch,\chi_{-4}\ch\). At \(x\asymp\sqrt p\)
their GRH PNT errors \(O(\sqrt x\log^2(px))\) are \(o(x)\).

\newpage
\section*{3. When the selected cycles capture every local candidate}
For \(q\in T\), choose either prime ideal \(\mathfrak q\) above \(q\),
put \(g=[\mathfrak q]\), and take any integer \(k\ge0\) with
\(Bq^k<\sqrt p\). If two of the sets
\[
 S,\quad gS,\quad\ldots,\quad g^k S
\]
intersect (at distinct indices), then \(q\) occurs in a prescribed cycle.
Indeed the intersection gives \(g^d C=C'\) with \(C,C'\in S\) and
\(1\le d\le k\). Choose a starting prime ideal \(\mathfrak s\) in \(C\).
The ideal \(\mathfrak s\mathfrak q^d\) in \(C'\) is primitive, since
\(s\ne q\) and the conjugate of neither prime ideal divides the product.
Its norm \(s q^d<\sqrt p\) makes it reduced. It must occur in the cycle
of each start in \(C'\).

Let \(k_B\) be the greatest integer \(k\ge0\) with \(B^{k+1}<\sqrt p\).
Writing \(W(B)\) for the distinct target primes in the full prescribed
cycles, we have the fixed-prime bounds
\begin{equation}
 |T|\,{\bf1}_{(k_B+1)|S|>h}\ \le\ |W(B)|\ \le\ |T|.
\end{equation}
The pigeonhole principle forces an intersection for each \(q\in T\).
For an individual \(q\) one may replace \(k_B\) by the larger admissible
\(k\) using its own size, and apply the same proof within any coset of
\(\langle g\rangle\) that meets \(S\). At \(B=\eta L^2\),
\[
 k_B+1\sim \frac{L}{4\log L}.
\]
Thus a fraction slightly greater than \(4\log L/L\) of the ordinary
classes suffices for complete detection, even if most classes are missed.
The number of starting roots alone gives no such lower bound on \(|S|/h\).

\textbf{Two individual checks.} For \(p=5711\) and \(B=0.26 L^2\),
the start \((71,5)\) gives \(5,47,26,67,10,95,\ldots\) and
\(95=5\cdot19\). The six starts from \(s=5,13,17\) occupy four of the
five ordinary classes. The sum in (1) before dividing by six is 10.
For a second direct check, the start \((72,17)\) has first neighbor 95.

For \(p=28{,}559\), \(B=0.11L^2=11.5788\ldots\), the only nonresidue
prime below \(B\) is 11 and the only starting prime is 5. Its two cycles
have the denominator lists
\[
 (5,67,10,199,137,134),\qquad(5,134,137,199,10,67).
\]
Both miss 11, as does the principal cycle with norms \(1,335,2,335\)
in cyclic order. Across all 15 ordinary classes there are 18 reduced
ideals with norm divisible by 11. The uniform class mean is \(18/15\).
The prescribed mean is zero. This finite example tests a coverage
inference and places no restriction on an eventual threshold for \(p\).
The accompanying Python check enumerates the reduced states and their exact
cycles at these two primes.

\newpage
\section*{4. What Dirichlet GRH supplies at the logarithmic scale}
Write \(c_*=\mathcal C(1)^{-2}<0.7615\) in the notation of Carneiro et
al. Their Theorem 1 guarantees a local nonresidue for fixed \(\eta>c_*\).
In the all-split case it is a member of \(T\). Under the coverage
condition (4), a prescribed cycle therefore succeeds at this scale.
There is also an all-prime guarantee for the direct test alone:
\begin{equation}
 \limsup_{\substack{p\to\infty\\p\equiv23\ (24)}}\
 \frac{\min\{s:\ s\text{ prime},\ s\equiv1\bmod4,\ \ch(s)=-1\}}{L^2}
 \ \le\ 4c_* <3.046.
\end{equation}
This is a specialization of the explicit-formula argument. The factor four
comes from restricting the prime sum to one of the two odd classes and then
squaring. Here are the details, including the dependence on the test.
\cmqr.

Fix \(F\) in their test class with \(\widehat F\) real and smooth of compact
support. Put \(x=\eta L^2,\ \Delta=\log(x)/(2\pi)\), and
\[
 D(F)=\int_{-\infty}^0\widehat F(t)e^{\pi t}\,dt
       -\int_0^\infty|\widehat F(t)|e^{\pi t}\,dt,\qquad
 \widehat h(t)=\tfrac12\{\widehat F(t-\Delta)+\widehat F(-t-\Delta)\}.
\]
For \(\xi=\ch,\ch\chi_{-4}\) of respective conductors \(f=p,4p\),
their explicit formula and zero sum bound give
\[
 \frac1\pi\sum_n\frac{\Lambda(n)\xi(n)}{\sqrt n}\,
       \widehat h\left(\frac{\log n}{2\pi}\right)
 \le\frac{\log f}{2\pi}\|F\|_1+o_F(L).
\]
Here \(\widehat h(0)=0\) eventually and each gamma integral is \(O_F(1)\).
Average the two formulas.
At odd primes the coefficient is \(\ch(s)\) on \(1\bmod4\) and zero on
\(3\bmod4\). The contribution from 2 and all higher prime powers is
\(O_F(\log^2 x)=o(L)\), since \(n\asymp_F x\) on the positive-frequency
support.

If the direct test misses at \(x\), insert \(+1\) on the selected primes
below \(x\) and bound the signed term from below by minus the absolute
value of its weight above \(x\). Partial summation with the PNT modulo 4
gives
\[
 \tfrac12\sqrt x\,D(F)+o_{F,\eta}(L)
 \ \le\ \frac{L}{2\pi}\|F\|_1+o_F(L).
\]
Indeed on setting \(n=xe^{2\pi t}\), the unrestricted prime measure
\(\pi^{-1}dn\,n^{-1/2}\widehat h(\log n/(2\pi))\) becomes
\(\sqrt x\,\widehat F(t)e^{\pi t}dt\). The progression contributes density
\(1/2\). The conductor bound above uses
\((\log p+\log(4p))/2=L+O(1)\). Infimizing
\(\|F\|_1/(2\pi D(F))\) over tests with \(D(F)>0\) gives
\(\mathcal C(1)^{-1}\) by their Theorem 4.
This proves (5). With all primes in place of \(1\bmod4\), the density factor
is one and the same calculation gives the local threshold \(c_*\).
In either calculation, a prime with the wrong sign below \(x\) can reduce
the lower bound by at most \(O_F(\log x/\sqrt x)\). Thus for a fixed
\(\eta\) strictly above the respective threshold there are
\(\gg_\eta x/\log x\) local witnesses, using a fixed test with strict margin.

\textbf{Consequence.} In the all-split case (4) transfers this last bound
to the cycles if the class coverage condition holds and \(\eta>c_*\).
At arbitrary small fixed \(\eta\), the missing inequality is
\(\sum_C w_C\nu_C>0\). Even with complete class coverage it would require
a split \(3\bmod4\) prime below \(\eta L^2\). Neither (2) nor (4) forces
one there. The pass gives no improvement to the growth rate of \(n_p\).
\end{document}
